\documentclass[10pt,a4paper]{amsart}
\usepackage[margin=19mm]{geometry}
\usepackage{amsmath,amssymb,mathtools}
\usepackage[scr=rsfs,cal=euler]{mathalfa}
\usepackage{xfrac}
\usepackage[unicode,hidelinks,bookmarks=false]{hyperref}
\newcommand{\ie}{i.e.\ }
\newcommand{\cf}{cf.\ }
\newcommand{\resp}{respectively\ }
\newcommand{\Subsection}{\subsection}
\providecommand{\nobreakdash}{\nobreak\hbox{-}}
\setlength{\emergencystretch}{2em}
\multlinegap0pt
\usepackage{amssymb}
\usepackage{tikz}
\newtheorem{lemma}{Lemma}
\title{Complemented zero-divisor graph of posets}
\author{Anagha Khiste, Ganesh Tarte, Vinayak Joshi}
\date{Source: arXiv:2604.15498v1 (CC BY 4.0)}
\begin{document}
\maketitle

\noindent\emph{Source and licence.} Complemented zero-divisor graph of posets. Version arXiv:2604.15498v1 (CC BY 4.0), distributed by arXiv under the Creative Commons Attribution 4.0 International licence, http://creativecommons.org/licenses/by/4.0/, as recorded in the arXiv licence metadata for this exact version and shown on the abstract page https://arxiv.org/abs/2604.15498v1. Full licence notice: \texttt{licenses/arxiv-2604.15498-CC-BY-4.0.txt}.\par\smallskip

\noindent\textit{Editorial scope.} Counted unit: the article's lemma 2.9 (source0.tex lines 148--150). The article's complete original proof of it is delivered with it (source0.tex lines 151--154, one \texttt{proof} environment of 4 lines). No mathematics is written, repaired or re-derived: the statement and the proof are the article's own lines, in the article's own notation, and no retained line is substituted or re-flowed.  No further source-local context is needed: every cross-reference of every retained line resolves inside the two delivered spans.  Nothing was omitted from any retained span, so the package carries no omission record.  No retained line contains a citation, so the wrapper carries no bibliography and no bibliography entry.  No retained line contains a figure environment, an image-inclusion command or drawing code, so no image is delivered and the package declares no asset dependency.  Editorial formatting only, in addition: an AMS \texttt{amsart} preamble that restates the article's own declarations and macros used by the retained lines, namely the environments \texttt{lemma} on separate counters, together with the standard packages those lines need.  The retained environments print in this document as Lemma 1 in order of appearance, whereas the article prints the same items with the numbering of its own full document.  The complete native source of the article as distributed in the arXiv e-print for this exact version is bundled unchanged as \texttt{sources/198609/source0.tex}.
\begin{lemma}\label{2.9} 
	Let $Q$ be a poset in $\mathbb{P}^{\ell}_{MFP}$. Then $int V(x)=(Spec(Q)\setminus \overline{D(x)})$ for any $x\in Q$.
\end{lemma}

\begin{proof} 
	Let $x\in Q$ be any element. First, we prove that $int V(x)\subseteq (Spec(Q)\setminus \overline{D(x)})$. Let $P\in int V(x)$. By the definition of the interior operator, $P\in D(a)$ for some open set $D(a)\subseteq V(x)$. Hence $D(a)\cap D(x) = \emptyset$. This implies the existence of a neighborhood $D(a)$ of $P$ such that $D(a)\cap D(x)=\emptyset $. Therefore, by the definition of the closure operator $P\notin \overline{D(x)}$, i.e., $P\in Spec(Q)\setminus \overline{D(x)}$. Hence $int V(x) \subseteq (Spec(Q)\setminus \overline{D(x)})$. Now, let $P\in Spec(Q)\setminus \overline{D(x)}$. Thus, there exists a neighborhood $D(a)$ of $P$ such that $D(a)\cap D(x) = \emptyset $. We claim that $\{a,x\}^{\ell}=\{0\}$. Suppose $\{a,x\}^{\ell}\neq \{0\}$. Then there exists a non-zero element $t$ such that $t\in \{a, x\}^{\ell}$. Since $Q\in \mathbb{P}^{\ell}_{MFP}$, the intersection of all prime ideals is $\{0\}$. Hence, there exists a prime ideal $P_1$ such that $t\not \in P_1$. But then $\{a, x\}^{\ell}\not \subseteq P_1$, implies that $P_1\in D(a)\cap D(x)$.  This contradicts the fact that $D(a)\cap D(x)=\emptyset $. Thus $\{a, x\}^{\ell}=\{0\}$, implying $D(a)\subseteq V(x)$. This implies $P\in int V(x)$. Hence, $(Spec(Q)\setminus \overline{D(x)})\subseteq int V(x)$. Therefore, $int V(x)=(Spec(Q)\setminus \overline{D(x)}).$ 
	%Note that, $\overline{int V(x)}=\overline{(Spec(Q)\setminus \overline{D(x)})}$.
\end{proof}

\end{document}
